\documentclass{article}
\usepackage[utf8]{inputenc}
\usepackage[T2A]{fontenc}
\usepackage[russian]{babel}
\usepackage{amsmath, amssymb, amsthm, amsfonts}
\usepackage{geometry}
\geometry{a4paper, margin=2cm}

\theoremstyle{plain}
\newtheorem{theorem}{Теорема}[section]
\newtheorem{lemma}[theorem]{Лемма}
\newtheorem{corollary}[theorem]{Следствие}
\newtheorem{definition}[theorem]{Определение}
\newtheorem{example}[theorem]{Пример}
\newtheorem{remark}{Замечание}[section]

\title{Конспект ответов на экзаменационные вопросы \\ по методам оптимизации}
\author{}
\date{}

\begin{document}

\maketitle

\section{Математическое программирование. Типы экстремумов функций многих переменных, условия локального экстремума, метод множителей Лагранжа, их интерпретация. Основные понятия выпуклого программирования. Седловые точки. Функция Лагранжа.}

\subsection{Определение математического программирования}
\textbf{Математическое программирование} — раздел математики, изучающий методы решения задач нахождения экстремумов (минимума или максимума) функций на множествах, задаваемых ограничениями (равенствами и неравенствами). Общая задача математического программирования имеет вид:
\[
f(x) \to \min \;(\text{или } \max),\quad x \in X \subset \mathbb{R}^n,
\]
где $f(x)$ — целевая функция, $X$ — допустимое множество, определяемое системой ограничений.

В зависимости от свойств функций и ограничений выделяют:
\begin{itemize}
  \item \textbf{Линейное программирование} — целевая функция и все ограничения линейны.
  \item \textbf{Нелинейное программирование} — хотя бы одна из функций (целевая или ограничения) нелинейна.
  \item \textbf{Выпуклое программирование} — целевая функция выпукла, допустимое множество выпукло (задаётся выпуклыми неравенствами и линейными равенствами).
  \item \textbf{Квадратичное программирование} — целевая функция квадратична, ограничения линейны.
  \item \textbf{Дискретное программирование} — переменные принимают только дискретные (например, целые) значения.
\end{itemize}

\subsection{Типы экстремумов}
\begin{definition}
Точка $x^* \in X$ называется точкой \textbf{глобального минимума} функции $f$ на $X$, если $f(x^*) \le f(x)$ для всех $x \in X$.
\end{definition}
\begin{definition}
Точка $x^* \in X$ называется точкой \textbf{локального минимума}, если существует $\varepsilon > 0$ такое, что $f(x^*) \le f(x)$ для всех $x \in X \cap U_\varepsilon(x^*)$, где $U_\varepsilon(x^*)$ — $\varepsilon$-окрестность.
\end{definition}
Аналогично определяются глобальный и локальный максимумы. В зависимости от наличия ограничений различают \textbf{безусловный} экстремум (когда $X = \mathbb{R}^n$) и \textbf{условный} (когда $X$ задаётся ограничениями).

Точка экстремума может быть:
\begin{itemize}
  \item \textbf{Стационарной} — если $f$ дифференцируема в этой точке и $\nabla f(x^*) = 0$ (необходимое условие экстремума).
  \item \textbf{Граничной} — если $x^*$ лежит на границе допустимого множества $X$ и не является внутренней точкой; в таких точках экстремум может достигаться без обращения градиента в ноль.
\end{itemize}

\subsection{Необходимые и достаточные условия локального экстремума}
Пусть $f: \mathbb{R}^n \to \mathbb{R}$ непрерывно дифференцируема в окрестности $x^*$.

\begin{theorem}[Необходимое условие (Ферма)]
Если $x^*$ — точка локального экстремума, то $\nabla f(x^*) = 0$.
\end{theorem}
\begin{proof}
Рассмотрим произвольное направление $d \in \mathbb{R}^n$, $\|d\|=1$, и функцию одной переменной $\varphi(t) = f(x^* + t d)$. При $t=0$ функция $\varphi$ имеет локальный экстремум, следовательно, $\varphi'(0) = \langle \nabla f(x^*), d \rangle = 0$. Здесь $\langle \cdot, \cdot \rangle$ — скалярное произведение в $\mathbb{R}^n$: $\langle u,v \rangle = \sum_{i=1}^n u_i v_i$. Поскольку $d$ произвольный, $\nabla f(x^*) = 0$.
\end{proof}

Пусть $f$ дважды непрерывно дифференцируема в окрестности $x^*$ и $\nabla f(x^*) = 0$. \textbf{Матрица Гессе} (гессиан) функции $f$ в точке $x^*$ определяется как
\[
H(x^*) = \left( \frac{\partial^2 f}{\partial x_i \partial x_j}(x^*) \right)_{i,j=1}^n.
\]

\begin{theorem}[Достаточное условие]
Если квадратичная форма $\langle H(x^*)y, y \rangle$ (где $y \in \mathbb{R}^n$ — произвольный ненулевой вектор, а $\langle H y, y \rangle = \sum_{i,j=1}^n H_{ij} y_i y_j$) строго положительно определена, то $x^*$ — точка локального минимума. Если строго отрицательно определена, то $x^*$ — точка локального максимума. Если форма неопределена, то экстремума нет.
\end{theorem}
\begin{proof}
Из разложения Тейлора:
\[
f(x^*+h) = f(x^*) + \frac{1}{2}\langle H(x^*)h, h\rangle + o(\|h\|^2).
\]
При положительной определённости $H(x^*)$ существует $\lambda > 0$ такое, что $\langle H(x^*)h, h\rangle \ge \lambda \|h\|^2$ для всех $h$. Тогда для достаточно малых $h$ имеем $f(x^*+h) - f(x^*) \ge \frac{\lambda}{4}\|h\|^2 > 0$, что означает локальный минимум. Аналогично для отрицательной определённости. Если форма неопределена, найдутся направления, по которым приращение положительно и отрицательно, значит, экстремума нет.
\end{proof}

Для проверки знакоопределённости применяется критерий Сильвестра: форма положительно определена $\iff$ все главные миноры матрицы Гессе положительны.

\subsection{Метод множителей Лагранжа}
Рассмотрим задачу с ограничениями-равенствами:
\[
f(x) \to \min,\quad \varphi_j(x)=0,\ j=1,\dots,m \quad (m \le n).
\]
Предполагаем, что функции непрерывно дифференцируемы и в точке экстремума $x^*$ матрица Якоби $\left(\frac{\partial \varphi_j}{\partial x_i}\right)$ имеет ранг $m$ (условие регулярности).

\begin{definition}
Функция Лагранжа: $L(x, \lambda) = f(x) + \sum_{j=1}^m \lambda_j \varphi_j(x)$, где $\lambda_j$ — множители Лагранжа.
\end{definition}

\begin{theorem}[Необходимые условия экстремума]
Если $x^*$ — точка экстремума задачи, то существуют такие $\lambda_1^*,\dots,\lambda_m^*$, что
\[
\frac{\partial L}{\partial x_i}(x^*,\lambda^*) = 0,\ i=1,\dots,n,\quad \frac{\partial L}{\partial \lambda_j}(x^*,\lambda^*) = \varphi_j(x^*) = 0.
\]
\end{theorem}
\begin{proof}[Идея доказательства]
При выполнении условия регулярности из теоремы о неявных функциях можно выразить $m$ переменных, например $x_1,\dots,x_m$, через остальные $x_{m+1},\dots,x_n$ в окрестности $x^*$. Тогда задача сводится к безусловной минимизации функции $\tilde{f}(x_{m+1},\dots,x_n) = f(x_1(x_{m+1},\dots,x_n),\dots,x_n)$. Условия стационарности для $\tilde{f}$ дают равенства, которые после применения правила дифференцирования сложной функции приводят к существованию множителей $\lambda_j^*$ и условиям $\frac{\partial L}{\partial x_i}=0$.
\end{proof}

\subsection{Экономическая интерпретация}
В задачах оптимизации ресурсов $\lambda_j$ интерпретируются как \textbf{теневые цены} (объективно обусловленные оценки). Они показывают, на сколько изменится оптимальное значение целевой функции при единичном увеличении ресурса $b_j$ (если ограничение имеет вид $\varphi_j(x) \le b_j$).

\subsection{Основные понятия выпуклого программирования}
\begin{definition}
Множество $X \subset \mathbb{R}^n$ выпукло, если $\forall x,y \in X$, $\forall \alpha \in [0,1]$: $\alpha x + (1-\alpha)y \in X$.
\end{definition}
\begin{definition}
Функция $f: X \to \mathbb{R}$ выпукла, если $\forall x,y \in X$, $\forall \alpha \in (0,1)$: $f(\alpha x + (1-\alpha)y) \le \alpha f(x) + (1-\alpha) f(y)$.
\end{definition}
Свойства: выпуклая функция непрерывна во внутренних точках; для дифференцируемой $f$ выпуклость эквивалентна $f(y) \ge f(x) + \langle \nabla f(x), y-x \rangle$; для дважды дифференцируемой — \textbf{положительной полуопределённости гессиана} (т.е. $\langle H(x) y, y \rangle \ge 0$ для всех $y$). Множество $\{x \in X \mid f(x) \le \lambda\}$ выпукло.

\subsection{Седловые точки и функция Лагранжа}
Рассматривается задача выпуклого программирования:
\[
f(x) \to \min,\quad \varphi_j(x) \le 0,\ j=1,\dots,m,
\]
где $f$ и $\varphi_j$ выпуклы.
Функция Лагранжа: $L(x, \lambda) = f(x) + \sum_{j=1}^m \lambda_j \varphi_j(x)$, $\lambda_j \ge 0$.

\begin{definition}
Пара $(x^*, \lambda^*)$, $\lambda^* \ge 0$, называется \textbf{седловой точкой}, если
\[
L(x^*, \lambda) \le L(x^*, \lambda^*) \le L(x, \lambda^*) \quad \forall x,\ \forall \lambda \ge 0.
\]
\end{definition}
\textbf{Прямая задача} — исходная задача минимизации $f(x)$ при ограничениях $\varphi_j(x)\le 0$. \textbf{Двойственная задача} — задача максимизации $\Theta(\lambda)=\inf_{x} L(x,\lambda)$ при $\lambda\ge 0$. Седловая точка даёт оптимальное решение прямой и двойственной задач: $x^*$ — решение прямой, $\lambda^*$ — решение двойственной.

\section{Теорема Куна – Таккера и ее геометрическая интерпретация. Современные методы градиентной оптимизации.}

\subsection{Теорема Куна – Таккера}
Рассмотрим задачу выпуклого программирования в дифференцируемой форме. Пусть $x^*$ — допустимая точка (т.е. $x^*$ удовлетворяет всем ограничениям), $I(x^*) = \{j \mid \varphi_j(x^*) = 0\}$.

\begin{theorem}[Куна – Таккера]
Если выполнено \textbf{условие регулярности Слейтера} (существует точка $x$ такая, что $\varphi_j(x) < 0$ для всех $j$), то $x^*$ — оптимальное решение тогда и только тогда, когда существуют такие $\lambda_j^* \ge 0$, $j \in I(x^*)$, что
\[
-\nabla f(x^*) = \sum_{j \in I(x^*)} \lambda_j^* \nabla \varphi_j(x^*). \tag{1}
\]
Для задачи с линейными ограничениями условие Слейтера не требуется.
\end{theorem}

\textit{Пояснение к доказательству (необходимость).} 
\textbf{Теория двойственности} — раздел выпуклого анализа, устанавливающий связь между исходной (прямой) задачей оптимизации и двойственной к ней. Для задачи $f(x)\to\min,\ \varphi_j(x)\le0$ двойственная задача имеет вид $\Theta(\lambda)=\inf_x L(x,\lambda)\to\max_{\lambda\ge0}$. Сильная двойственность означает, что оптимальные значения прямой и двойственной задач совпадают.

\textbf{Конус из градиентов} — множество всех неотрицательных линейных комбинаций градиентов активных ограничений:
\[
K = \left\{ \sum_{j\in I(x^*)} \lambda_j \nabla \varphi_j(x^*) \mid \lambda_j \ge 0 \right\}.
\]
Условие (1) означает $-\nabla f(x^*) \in K$, т.е. антиградиент целевой функции лежит в этом конусе. Геометрически это исключает существование направления спуска, не выводящего за пределы допустимого множества.

\textit{Доказательство необходимости:} Пусть $x^*$ — оптимальное решение. Предположим, что $-\nabla f(x^*)$ не принадлежит конусу $K$. Тогда по теореме отделимости существует вектор $d$ такой, что $\langle -\nabla f(x^*), d \rangle > 0$ и $\langle \nabla \varphi_j(x^*), d \rangle \le 0$ для всех $j\in I(x^*)$. Первое неравенство означает $\langle \nabla f(x^*), d \rangle < 0$, т.е. $d$ — направление убывания $f$. Второе — что для малых шагов $x^*+td$ остаётся допустимой (ограничения не нарушаются). Это противоречит оптимальности $x^*$. Следовательно, $-\nabla f(x^*)\in K$, что и даёт (1). Достаточность следует из выпуклости функций и условий дополняющей нежёсткости.

\subsection{Современные методы градиентной оптимизации}
\subsubsection{Градиентный спуск}
Итерационная формула: $x^{k+1} = x^k - h_k \nabla f(x^k)$.

\textbf{Определение.} Функция $f$ имеет $L$-липшицев градиент, если для всех $x,y$ выполнено
\[
\|\nabla f(y) - \nabla f(x)\|_2 \le L \|y - x\|_2,
\]
где $\|v\|_2 = \sqrt{v_1^2+\dots+v_n^2}$ — евклидова норма. Константа $L>0$ называется \textbf{константой Липшица} градиента. Она характеризует максимальную скорость изменения градиента: чем меньше $L$, тем «более гладкая» функция. Это свойство гарантирует квадратичную оценку сверху:
\[
f(y) \le f(x) + \langle \nabla f(x), y-x \rangle + \frac{L}{2}\|y-x\|_2^2.
\]

\textbf{Определение.} Функция $f$ называется $\mu$-сильно выпуклой, если существует $\mu>0$ (константа сильной выпуклости) такое, что для всех $x,y$
\[
f(y) \ge f(x) + \langle \nabla f(x), y-x \rangle + \frac{\mu}{2}\|y-x\|_2^2.
\]
Параметр $\mu$ определяет, насколько сильно функция «изогнута» вниз. Чем больше $\mu$, тем быстрее гарантированно сходится градиентный спуск. Для квадратичной функции $f(x)=\frac12 x^T A x$ константа $\mu$ равна минимальному собственному числу матрицы $A$. Это свойство обеспечивает линейную сходимость методов.

Для выпуклых функций с $L$-липшицевым градиентом выбор $h = 1/L$ даёт
\[
f(\bar{x}^N) - f(x^*) \le \frac{L R^2}{2N},\quad R = \|x^0 - x^*\|_2,\quad \bar{x}^N = \frac{1}{N}\sum_{k=1}^N x^k.
\]
Для $\mu$-сильно выпуклых функций: $f(x^{N+1})-f(x^*) \le (1-\mu/L)^N (f(x^0)-f(x^*))$.

\subsubsection{Ускоренный градиентный метод (Нестерова)}
Двухшаговая схема:
\[
y^{k+1} = x^k + \frac{k}{k+3}(x^k - x^{k-1}),\quad x^{k+1} = y^{k+1} - \frac{1}{L}\nabla f(y^{k+1}).
\]
Оценка: $f(x^N) - f(x^*) \le \frac{2LR^2}{N^2}$, что оптимально для гладких выпуклых задач.

\subsubsection{Стохастический градиентный спуск}
Применяется для $f(x) = \mathbb{E}_{\xi}[f(x,\xi)]$, когда точный градиент недоступен, но можно получить его несмещённую оценку. На итерации используется $\nabla f(x^k, \xi^k)$, где $\xi^k$ — случайная величина (реализация случайного вектора). Оценки сходимости: $\mathbb{E}[f(\bar{x}^N)-f(x^*)] = O(1/\sqrt{N})$. При этом дисперсия оценки влияет на константу.

\subsection{Сравнение идей градиентных методов}

\textbf{Обычный градиентный спуск (2.2.1)} — самый простой метод. Идея: на каждой итерации двигаться в сторону наискорейшего убывания функции (антиградиент) с фиксированным или адаптивным шагом. Этот метод похож на движение шарика, который всегда скатывается по самому крутому склону. Он надёжен, но может быть медленным в «оврагах» (когда линии уровня сильно вытянуты).

\textbf{Ускоренный градиентный спуск (Нестерова, 2.2.2)} — добавляет «импульс» (инерцию). Идея: использовать не только текущий антиградиент, но и направление с предыдущего шага. Это похоже на движение тяжёлого шарика, который не останавливается на каждом локальном спуске, а продолжает катиться, накапливая скорость. Благодаря этому метод «проскакивает» мелкие неровности и быстрее сходится к оптимуму, особенно в оврагах. Теоретическая оценка сходимости $O(1/N^2)$ вместо $O(1/N)$ у обычного спуска.

\textbf{Стохастический градиентный спуск (2.2.3)} — используется, когда точный градиент вычислить дорого или невозможно (например, функция задана как сумма по большому числу слагаемых). Идея: вместо точного градиента брать его случайную оценку (например, градиент одного случайно выбранного слагаемого). Это похоже на движение в условиях шума: каждый шаг не точно направлен к минимуму, но в среднем направление верное. Цена — меньшая скорость сходимости ($O(1/\sqrt{N})$), но зато каждый шаг очень дёшев. Этот метод лежит в основе обучения нейронных сетей.

\textbf{Ключевое отличие:}
\begin{itemize}
\item Обычный спуск — честно идёт по антиградиенту.
\item Ускоренный спуск — добавляет инерцию, чтобы быстрее проходить овраги.
\item Стохастический спуск — жертвует точностью направления ради дешевизны шага.
\end{itemize}

\section{Формулировка задачи линейного программирования (ЛП). Понятия опорного плана и базиса, вырожденность и невырожденность задач ЛП, основные принципы симплекс-метода. Основные теоремы ЛП.}

\subsection{Постановка задачи ЛП}
Каноническая форма:
\[
w(x) = c^T x \to \min,\quad Ax = b,\ x \ge 0,
\]
где $A \in \mathbb{R}^{m \times n}$, $\operatorname{rank} A = m$, $b \ge 0$. 
Вектор $c = (c_1,\dots,c_n)^T$ — это \textbf{вектор коэффициентов целевой функции} (стоимостей, весов). Целевая функция записывается как $w(x)=c^T x = \sum_{j=1}^n c_j x_j$. Например, если $x_j$ — количество выпускаемой продукции $j$, а $c_j$ — прибыль от единицы продукции, то $w(x)$ — общая прибыль (максимизируется) или затраты (минимизируется). Переменные $x_j \ge 0$ называются неотрицательными. Допустимое множество $X = \{x \mid Ax = b,\ x \ge 0\}$ — выпуклый многогранник (возможно, неограниченный).

\subsection{Базис, опорный план, вырожденность}
Пусть $B = (A_{\sigma(1)},\dots,A_{\sigma(m)})$ — набор из $m$ линейно независимых столбцов матрицы $A$ (базис). Соответствующие переменные $x_B$ — \textbf{базисные}, остальные $x_N$ — \textbf{свободные}. Базисное решение: $x_B = B^{-1}b$, $x_N = 0$. Здесь $b = (b_1,\dots,b_m)^T$ — \textbf{вектор правых частей ограничений} (ресурсов). Если $x_B \ge 0$, то оно называется \textbf{опорным планом} (базисным допустимым решением).

\begin{definition}
Опорный план называется \textbf{вырожденным}, если среди его базисных переменных есть нулевые (т.е. $|\{j: x_j > 0\}| < m$). В противном случае план называется \textbf{невырожденным}.
\end{definition}
Вырожденность возникает, когда в системе ограничений есть избыточные связи или когда базисное решение находится на пересечении более чем $m$ гиперплоскостей. Наличие вырожденности может приводить к зацикливанию в симплекс-методе.

\subsection{Симплекс-метод}
Основная идея: переход от одного опорного плана к смежному (отличающемуся одной базисной переменной) с улучшением целевой функции.

Симплекс-таблица строится по базису. Входные данные: матрица $A$, векторы $b$, $c$, начальный допустимый базис (если его нет — используется метод искусственного базиса). Выходные данные: оптимальный план $x^*$ (если существует) или заключение о неограниченности/несовместности.

Обозначения в таблице:
\begin{itemize}
  \item $z_{00} = -c_B B^{-1} b$ — текущее значение целевой функции со знаком минус.
  \item $z_{0j} = c_j - c_B B^{-1} A_j$ — оценки (приращения) для небазисных переменных.
  \item $z_{i0}$ — значения базисных переменных.
  \item $z_{ij}$ — коэффициенты разложения $j$-го столбца по базису.
\end{itemize}

Критерий оптимальности: если все $z_{0j} \ge 0$, то текущий план оптимален.

Алгоритм:
\begin{enumerate}
\item Найти $s$ с $z_{0s} < 0$ (ведущий столбец). Если таких нет — конец (оптимум).
\item Если все $z_{is} \le 0$, то целевая функция неограничена снизу.
\item Иначе выбрать $r = \arg\min_{i: z_{is}>0} \frac{z_{i0}}{z_{is}}$ (ведущая строка). Элемент $z_{rs}$ — ведущий.
\item Преобразовать таблицу (элементарные преобразования Гаусса), перевести $x_s$ в базис, $x_{\sigma(r)}$ — в свободные.
\item Перейти к шагу 1.
\end{enumerate}

\subsection{Основные теоремы ЛП}
\begin{theorem}[Критерий разрешимости]
Задача ЛП разрешима (т.е. существует оптимальное решение) тогда и только тогда, когда допустимое множество непусто и целевая функция ограничена снизу на нём.
\end{theorem}
\begin{theorem}[Существование опорного плана]
Если допустимое множество непусто, то существует опорный план (базисное допустимое решение).
\end{theorem}
\begin{theorem}[Оптимальность на опорном плане]
Если задача разрешима, то существует оптимальный опорный план.
\end{theorem}
\begin{theorem}[Двойственности]
Прямая и двойственная задачи либо обе разрешимы (и тогда их оптимальные значения равны), либо обе неразрешимы (одна из-за несовместности, другая из-за неограниченности). Для оптимальных планов $x^*$, $y^*$ выполняются условия дополняющей нежёсткости: $y_i^*(a_i x^* - b_i) = 0$, $(c_j - y^*A_j)x_j^* = 0$.
\end{theorem}

\section{Потоки в сетях. Теорема Форда – Фалкерсона.}
\subsection{Основные определения}
\begin{definition}
Сеть — ориентированный граф $G=(V,E)$ с источником $s$ и стоком $t$. Каждой дуге $(u,v)$ приписана пропускная способность $c(u,v) \ge 0$. \textbf{Поток} $f$ — функция на дугах, удовлетворяющая:
\begin{enumerate}
\item $0 \le f(u,v) \le c(u,v)$;
\item $\sum_{v} f(u,v) = \sum_{v} f(v,u)$ для всех $u \neq s,t$ (сохранение потока).
\end{enumerate}
Величина потока $|f| = \sum_{v} f(s,v) - \sum_{v} f(v,s)$.
\end{definition}

\textbf{Разрез} $(S,T)$: $s \in S$, $t \in T$, $T = V \setminus S$. Пропускная способность разреза $C(S,T) = \sum_{u\in S, v\in T} c(u,v)$. Переменные: $f(u,v)$ — поток по дуге, $c(u,v)$ — пропускная способность.

\subsection{Теорема Форда – Фалкерсона}
\begin{theorem}[Форда – Фалкерсона]
Максимальная величина потока равна минимальной пропускной способности разреза.
\end{theorem}
\begin{proof}
\textit{1. Оценка сверху:} Для любого потока $f$ и любого разреза $(S,T)$ имеем $|f| = \sum_{u\in S, v\in T} f(u,v) - \sum_{u\in S, v\in T} f(v,u) \le \sum_{u\in S, v\in T} c(u,v) = C(S,T)$. Следовательно, максимальный поток не превышает минимальной пропускной способности разреза.

\textit{2. Построение потока, достигающего этой границы:} Алгоритм Форда – Фалкерсона строит увеличивающие пути в остаточной сети. Остаточная сеть содержит прямые дуги с остаточной пропускной способностью $c(u,v)-f(u,v)$ и обратные дуги с $f(u,v)$. Если увеличивающего пути из $s$ в $t$ в остаточной сети нет, то множество $S$ вершин, достижимых из $s$ по ненасыщенным дугам, не содержит $t$. Для этого разреза $(S,T)$ все прямые дуги из $S$ в $T$ насыщены ($f(u,v)=c(u,v)$), а обратные дуги имеют нулевой поток ($f(v,u)=0$). Следовательно, $|f| = \sum_{u\in S, v\in T} c(u,v) = C(S,T)$. Таким образом, поток достигает пропускной способности разреза, и по предыдущему пункту он максимален.

\textit{3. Конечность при целых пропускных способностях:} При целых $c(u,v)$ алгоритм увеличивает поток на целое положительное число на каждом шаге, поэтому за конечное число шагов будет достигнут максимум.
\end{proof}

\section{Динамическое программирование. Примеры задач, решаемых методом динамического программирования.}

\subsection{Общая постановка и принцип оптимальности}
Рассматривается $N$-шаговый процесс. \textbf{Состояние} на $k$-м шаге $x^{(k)}$ — это вектор параметров, полностью описывающий систему в данный момент. \textbf{Управление} $u^{(k)}$ — решение, принимаемое на $k$-м шаге, переводящее систему из состояния $x^{(k-1)}$ в $x^{(k)}$.

Уравнение состояний: $x^{(k)} = f^{(k)}(x^{(k-1)}, u^{(k)})$. \textbf{Целевая функция} (критерий эффективности) аддитивна: $J = \sum_{k=1}^N J_k(x^{(k-1)}, u^{(k)})$.

\begin{definition}[Принцип оптимальности Беллмана]
Для любого начального состояния оптимальная стратегия на последующих шагах не зависит от того, как система пришла в это состояние, и определяется только самим состоянием.
\end{definition}

Уравнения Беллмана (для максимизации):
\[
B_N(x^{(N-1)}) = \max_{u^{(N)}} J_N(x^{(N-1)}, u^{(N)}),
\]
\[
B_k(x^{(k-1)}) = \max_{u^{(k)}} \left\{ J_k(x^{(k-1)}, u^{(k)}) + B_{k+1}(x^{(k)}) \right\},\quad k = N-1,\dots,1.
\]
Результатом является $B_1(x^{(0)})$ — максимальное значение целевой функции за все $N$ шагов, и условно-оптимальные управления $u^{*(k)}(x^{(k-1)})$. Затем, зная $x^{(0)}$, последовательно находятся $x^{(1)},\dots,x^{(N)}$ и $u^{(1)},\dots,u^{(N)}$.

\subsection{Примеры задач}

\subsubsection{Задача о распределении средств между предприятиями}
Имеется $N$ предприятий, начальные средства $x^{(0)}$. Прибыль от вложения $u$ в $k$-е предприятие $J_k(u)$. Требуется распределить средства так, чтобы суммарная прибыль была максимальной.

Параметр состояния $x^{(k)}$ — остаток средств после $k$ предприятий. Уравнение состояний: $x^{(k)} = x^{(k-1)} - u^{(k)}$. Уравнения Беллмана решаются таблично, начиная с последнего предприятия. Пример из пособия Ларина: при $x^{(0)}=5$ и заданных $J_k(u)$ оптимальное распределение $(1,2,1,1)$ с суммарной прибылью $24$.

\subsubsection{Задача об оптимальном распределении ресурсов между отраслями на $N$ лет}
Две отрасли, начальные средства $x^{(0)}$. Возврат средств: $q_1(u)$ и $q_2(x-u)$. Прибыль $J_1(u)+J_2(x-u)$. Уравнение состояний: $x^{(k)} = q_1(u^{(k)}) + q_2(x^{(k-1)}-u^{(k)})$. Решение рекуррентных соотношений приводит к оптимальной стратегии инвестиций.

\subsubsection{Задача о замене оборудования}
Оборудование эксплуатируется $N$ лет. Затраты на содержание $r(t)$, ликвидная стоимость $g(t)$. Решение: в начале каждого года принимать решение — сохранить или заменить. Параметр состояния — возраст $t$. Уравнения Беллмана минимизируют суммарные затраты. В примере из пособия при $N=5$ оптимальная стратегия: заменить в начале 3-го года.

% ============= ДОБАВЛЕННЫЕ РАЗДЕЛЫ =============

\section*{Приложение А. Примеры решения задач}

\subsection*{Пример 1 (метод множителей Лагранжа)}
Найти экстремум функции $f(x,y)=x^2+y^2$ при условии $x+y=2$.

\textit{Решение.} Составляем функцию Лагранжа $L(x,y,\lambda)=x^2+y^2+\lambda(2-x-y)$. Условия стационарности:
\[
\frac{\partial L}{\partial x}=2x-\lambda=0,\quad \frac{\partial L}{\partial y}=2y-\lambda=0,\quad \frac{\partial L}{\partial \lambda}=2-x-y=0.
\]
Из первых двух $x=y$, подставляем в третье: $2-2x=0 \Rightarrow x=1$, $y=1$, $\lambda=2$. Гессиан $H=\begin{pmatrix}2&0\\0&2\end{pmatrix}$ положительно определён, поэтому $(1,1)$ — точка минимума, $f_{\min}=2$.

\subsection*{Пример 2 (симплекс-метод)}
Решить задачу ЛП: $w=-x_1-2x_2\to\min$, при $x_1+x_2\le 4$, $x_1+2x_2\le 6$, $x_1,x_2\ge0$.

\textit{Решение.} Приводим к канонической форме: $x_1+x_2+x_3=4$, $x_1+2x_2+x_4=6$, $x_i\ge0$, $w=-x_1-2x_2\to\min$. Начальный базис $B=(A_3,A_4)$: $x_3=4$, $x_4=6$, $x_1=x_2=0$, $w=0$. Симплекс-таблица:
\[
\begin{array}{c|cccc|c}
 & x_1 & x_2 & x_3 & x_4 & \\ \hline
-w & 0 & -1 & -2 & 0 & 0 \\
x_3 & 4 & 1 & 1 & 1 & 0 \\
x_4 & 6 & 1 & 2 & 0 & 1
\end{array}
\]
Ведущий столбец $x_2$ ($z_{02}=-2<0$). Отношения: $\frac{4}{1}=4$, $\frac{6}{2}=3$ → ведущая строка $x_4$, ведущий элемент $2$. После преобразования получаем новую таблицу:
\[
\begin{array}{c|cccc|c}
 & x_1 & x_2 & x_3 & x_4 & \\ \hline
-w & 6 & 0 & 0 & 1 & 1 \\
x_3 & 1 & 1/2 & 0 & 1 & -1/2 \\
x_2 & 3 & 1/2 & 1 & 0 & 1/2
\end{array}
\]
Все $z_{0j}\ge0$ → оптимальное решение: $x_1=0$, $x_2=3$, $w_{\min}=-6$.

\subsection*{Пример 3 (теорема Форда – Фалкерсона)}
Сеть: $s\to a$ (3), $s\to b$ (2), $a\to t$ (2), $b\to t$ (3), $a\to b$ (1). Найти максимальный поток.

\textit{Решение.} Начальный нулевой поток. Увеличивающий путь $s\to a\to t$: увеличиваем на 2 (насыщается дуга $a\to t$). Далее путь $s\to b\to t$: увеличиваем на 2 (насыщается $s\to b$). Затем путь $s\to a\to b\to t$: можно добавить 1 (по $a\to b$). Итоговый поток: $f(s,a)=3$, $f(s,b)=2$, $f(a,t)=2$, $f(b,t)=3$, $f(a,b)=1$. Величина потока $=5$. Минимальный разрез: $S=\{s,a,b\}$, $T=\{t\}$, пропускная способность $c(a,t)+c(b,t)=2+3=5$. Теорема подтверждается.

\subsection*{Пример 4 (динамическое программирование)}
Задача о распределении средств: $x^{(0)}=4$, 3 предприятия. Прибыль: $J_1(1)=5$, $J_1(2)=8$, $J_1(3)=10$; $J_2(1)=6$, $J_2(2)=9$, $J_2(3)=11$; $J_3(1)=4$, $J_3(2)=7$, $J_3(3)=12$. Найти оптимальное распределение.

\textit{Решение.} Шаг 3 (последнее предприятие): $B_3(x^{(2)})=\max_{0\le u\le x^{(2)}} J_3(u)$. $B_3(0)=0$, $B_3(1)=4$, $B_3(2)=7$, $B_3(3)=12$, $B_3(4)=12$ (лучше вложить 3). Шаг 2: $B_2(x^{(1)})=\max_{0\le u\le x^{(1)}} [J_2(u)+B_3(x^{(1)}-u)]$. Вычисляем: при $x^{(1)}=4$: $u=0$: $0+B_3(4)=12$; $u=1$: $6+B_3(3)=6+12=18$; $u=2$: $9+B_3(2)=9+7=16$; $u=3$: $11+B_3(1)=11+4=15$; $u=4$: нет. Максимум 18 при $u=1$. Шаг 1: $B_1(x^{(0)})=\max_{u}[J_1(u)+B_2(4-u)]$. $u=0$: $0+B_2(4)=18$; $u=1$: $5+B_2(3)=5+?$ $B_2(3)$: $u=0$: $0+B_3(3)=12$; $u=1$: $6+B_3(2)=6+7=13$; $u=2$: $9+B_3(1)=9+4=13$; $u=3$: $11+B_3(0)=11$ → $B_2(3)=13$. Итого $5+13=18$; $u=2$: $8+B_2(2)=8+?$ $B_2(2)$: $u=0$: $0+B_3(2)=7$; $u=1$: $6+B_3(1)=6+4=10$; $u=2$: $9+B_3(0)=9$ → $B_2(2)=10$, сумма $18$; $u=3$: $10+B_2(1)=10+?$ $B_2(1)$: $u=0$: $0+B_3(1)=4$; $u=1$: $6+B_3(0)=6$ → $B_2(1)=6$, сумма $16$; $u=4$: нет. Оптимум 18 достигается при $u=1$ (первому предприятию 1) или $u=2$ (2). Выбираем, например, $u^{(1)}=1$, тогда $x^{(1)}=3$, $u^{(2)}=1$ (из $B_2(3)$ при $u=1$), $x^{(2)}=2$, $u^{(3)}=2$ (из $B_3(2)$). Распределение $(1,1,2)$, прибыль 18.

\section*{Приложение Б. Иллюстрации (описание)}

Ниже приведены краткие описания ключевых рисунков, которые полезно иметь в конспекте.

\begin{enumerate}
  \item \textbf{Геометрическая интерпретация градиентного спуска:} Параболоид, касающийся графика функции в текущей точке и мажорирующий его; точка минимума параболоида даёт следующее приближение.
  \item \textbf{Симплекс-метод:} Таблица с базисными и свободными переменными, ведущий столбец и ведущая строка, преобразование по методу Гаусса.
  \item \textbf{Сеть и разрез:} Ориентированный граф с источником $s$ и стоком $t$, на дугах указаны пропускные способности; разрез — множество дуг, удаление которых разъединяет $s$ и $t$.
  \item \textbf{Динамическое программирование:} Схема переходов между состояниями (многошаговый процесс), где каждый шаг характеризуется состоянием и управлением; стрелки показывают возможные переходы.
\end{enumerate}

\section*{Приложение В. Сравнительная таблица методов оптимизации}

\begin{center}
\renewcommand{\arraystretch}{1.3}
\begin{tabular}{|p{2.8cm}|p{2.2cm}|p{2.5cm}|p{2.5cm}|p{2.5cm}|}
\hline
\textbf{Метод} & \textbf{Область применимости} & \textbf{Входные данные} & \textbf{Сложность (число итераций)} & \textbf{Типичные приложения} \\
\hline
Градиентный спуск & Гладкие (выпуклые) функции без ограничений & $f(x)$, $\nabla f(x)$, $x^0$ & $O(1/\varepsilon)$ (выпуклые) & Обучение нейронных сетей, регрессия \\
\hline
Симплекс-метод & Линейное программирование & $A$, $b$, $c$, начальный базис & $O(2^m)$ в худшем, $O(m^3)$ в среднем & Планирование производства, транспорт \\
\hline
Форда–Фалкерсона & Максимальный поток в сети & Граф, пропускные способности, $s$, $t$ & $O(|E|\cdot f_{\max})$ & Сети связи, логистика \\
\hline
Динамическое программирование & Многошаговые процессы & Рекуррентные соотношения, начальное состояние & $O(N\cdot |X|^2)$ (зависит от дискретизации) & Управление запасами, замена оборудования \\
\hline
\end{tabular}
\end{center}

\section*{Приложение Г. Практические замечания по выбору метода}

\begin{enumerate}
  \item \textbf{Если задача малой размерности (до 10 переменных) и функции гладкие:} можно использовать аналитические методы (множители Лагранжа) или градиентные методы с контролем точности.
  \item \textbf{Если задача линейная и имеет много переменных:} предпочтителен симплекс-метод или методы внутренней точки; симплекс-метод эффективен при числе ограничений до нескольких тысяч.
  \item \textbf{Если задача нелинейная, но выпуклая:} хороши градиентные методы (особенно ускоренные) и метод Ньютона (если есть гессиан). При наличии ограничений — метод проекции градиента или метод условного градиента.
  \item \textbf{Если задача о потоках в сети:} алгоритм Форда–Фалкерсона или его модификации (Эдмондса–Карпа) дают гарантированный результат за полиномиальное время.
  \item \textbf{Если процесс многошаговый и может быть разбит на этапы:} динамическое программирование — лучший выбор, особенно когда состояния дискретны.
  \item \textbf{При большом количестве переменных ($n > 10^4$) и отсутствии точного градиента:} используют стохастический градиентный спуск или его варианты (Adam, RMSprop).
\end{enumerate}

\section*{Приложение Д. Глоссарий}

\begin{itemize}
  \item \textbf{Градиент} — вектор частных производных функции: $\nabla f = \left(\frac{\partial f}{\partial x_1},\dots,\frac{\partial f}{\partial x_n}\right)^T$.
  \item \textbf{Матрица Гессе (гессиан)} — матрица вторых частных производных: $H_{ij}=\frac{\partial^2 f}{\partial x_i\partial x_j}$.
  \item \textbf{Опорный план} — базисное допустимое решение задачи ЛП (все компоненты неотрицательны).
  \item \textbf{Вырожденность} — наличие нулевых базисных переменных в опорном плане.
  \item \textbf{Седловая точка} — точка $(x^*,\lambda^*)$, для которой $L(x^*,\lambda)\le L(x^*,\lambda^*)\le L(x,\lambda^*)$.
  \item \textbf{L-липшицев градиент} — свойство $\|\nabla f(y)-\nabla f(x)\|\le L\|y-x\|$ (евклидова норма).
  \item \textbf{$\mu$-сильная выпуклость} — свойство $f(y)\ge f(x)+\langle\nabla f(x),y-x\rangle+\frac{\mu}{2}\|y-x\|^2$.
  \item \textbf{Остаточная сеть} — сеть, в которой для каждой дуги $(u,v)$ добавлена обратная дуга с пропускной способностью, равной текущему потоку; используется в алгоритме Форда–Фалкерсона.
  \item \textbf{Увеличивающий путь} — путь из источника в сток в остаточной сети, по которому можно увеличить поток.
  \item \textbf{Принцип оптимальности Беллмана} — основа динамического программирования: оптимальное управление на оставшихся шагах зависит только от текущего состояния.
\end{itemize}

\end{document}
